\chapter{Data, backups and releases}
\label{ch:backups}

\section{Data}\label{data}

Git isn\textquotesingle t built for large files, so keep data out of it, and record what you used.

\begin{itemize}
\tightlist
\item
  \textbf{Never overwrite.} Write each new run to a dated folder, such as \texttt{results\_2026-10-31/}.
\item
  \textbf{Record checksums}, which are fingerprints that change if a file changes: \texttt{sha256sum\ *.csv\ \textgreater{}\ MANIFEST.sha256}, and check later with \texttt{sha256sum\ -c\ MANIFEST.sha256}.
\item
  \textbf{Note the source} of each dataset: its address, date and licence.
\end{itemize}

\textbf{What you should see when nothing has changed}

\begin{lstlisting}[language=text]
$ sha256sum -c MANIFEST.sha256
a.csv: OK
b.csv: OK
\end{lstlisting}

\textbf{What you should see after someone edits b.csv}

\begin{lstlisting}[language=text]
$ sha256sum -c MANIFEST.sha256
a.csv: OK
b.csv: FAILED
sha256sum: WARNING: 1 computed checksum did NOT match
\end{lstlisting}

\section{Backups}\label{backups}

Git is not a backup until a copy exists elsewhere, and a private repository on one laptop can be lost with that laptop. The common rule is \textbf{3-2-1}: three copies, on two kinds of storage, with one somewhere else.

Syncing services such as OneDrive aren\textquotesingle t enough on their own, because they also sync deletions and corruption. And \textbf{a backup you\textquotesingle ve never restored is only a hope}, so test one:

\begin{lstlisting}[language=bash]
git bundle create project.bundle --all     # one file holding a repository's whole history
git clone project.bundle restore-test      # prove that it comes back
\end{lstlisting}

\texttt{git\ bundle\ create} prints nothing when it runs inside a script, as here, and shows progress when you type it yourself. Then read the restored history to check it:

\textbf{What you should see}

\begin{lstlisting}[language=text]
$ git clone project.bundle restore-test
Cloning into 'restore-test'...

$ git -C restore-test log --oneline -3
4912949 Colleague edit
b889be1 Merge left, keeping 14
a0a50db Main change
\end{lstlisting}

Databases and uploaded files don\textquotesingle t live in Git, so they need their own backups, including an encrypted copy off-site.

\section{Releases}\label{releases}

Whenever something leaves your hands, freeze a copy.

\begin{enumerate}
\def\labelenumi{\arabic{enumi}.}
\tightlist
\item
  Commit everything, so that the source is exact.
\item
  Run all the checks.
\item
  Build from clean.
\item
  Put the outputs in a dated, read-only folder, with checksums.
\item
  Tag the commit.
\end{enumerate}

Then, for anything a reader sends back, you can find precisely what they had. If it needs correcting, make a new release and log the reason, without editing the old one.

\section{Secrets}\label{secrets}

Passwords and keys never go in a repository, because history keeps them even after deletion. Use environment variables or a secrets store. If one gets committed, treat it as exposed: change it first, then clean up.

